\section{Ethics and Broader Impact}

This work studies backdoored adapters that can induce harmful compliance under a trigger.
The purpose is defensive: to understand whether post-hoc adapter transformations can reduce malicious behavior when retraining the base model is unavailable.
The experiments therefore call for careful disclosure.
Harmful prompts, trigger strings, and target completions are handled under controlled access or anonymized where disclosure would materially aid misuse.
At the same time, the public release supports independent verification with sanitized evaluation scripts, aggregate metrics, component gates, a dataset card, annotation guidelines, sanitized examples, training-metadata status, and a synthetic trigger setting that exercises the same protocol.

The main deployment risk is over-refusal.
A defense that removes attack success by refusing all triggered inputs may reduce immediate harm but can create a denial-of-service behavior or degrade legitimate benign interactions.
For this reason, we report triggered-benign and benign refusal separately from triggered harmful ASR: Qwen4B has usable frontier points, but its most aggressive high-refusal endpoints are not the preferred deployment choices.

The broader implication is that adapter compression should not be treated as a purely efficiency-oriented step.
Compression can preserve, attenuate, or redirect unsafe behaviors.
Practitioners who load third-party adapters should evaluate safety after compression or quantization, not only before it.
